\section{Related literature}

The research technology belongs to the tradition of endogenous and
semi-endogenous growth models in which resources devoted to ideas determine the
evolution of productivity \citep{romer1990,aghionhowitt1992,jones1995}. The
population term matters for the same reason it matters in \citet{jones1995}: when
human researchers remain essential, growth in the research workforce can anchor
long-run innovation. I ask when that anchor disappears because
capability-augmented research services can substitute for human researchers. I use
a homothetic CES research technology to separate substitution between research
inputs from returns to their current joint scale. The latter role of \(\eta\) is analogous
to the intratemporal research-scale parameter in \citet{jones1995}; the baseline
fixes the separate stock-of-ideas exponent at zero, corresponding to \(b=1\) in
the extension discussed below.

The closest paper for the feedback architecture is \citet{davidsonetal2026}. They model innovation as a network
of technological and economic feedbacks, classify its growth regimes using the
spectral radius of a nonnegative feedback matrix, and combine nonrival AI
capability with rival compute. Their application distinguishes software, hardware
quality, and general productivity, but fixes saving and resource-allocation shares
and treats task automation as exogenous, including an exogenous path in its CES
extension. I collapse that technology network into one
capability stock and place its two feedback loops inside a decentralized Ramsey
equilibrium with an integrated monopolist, endogenous inference and research
compute, endogenous labor allocation, two CES substitution margins, and factor
prices. The mapping is not literal: in my formulation, \(U\) and \(M\) are flow
resource uses, not separate accumulable hardware stocks. The unit-elasticity
denominator in my model is therefore a scalar analogue of their
feedback criterion, not an application of their general network theorem.

The large-scale curvature restriction derived below answers a different
question from \citet{davidsonetal2026}'s feedback criterion. Their matrix
criterion classifies a coupled system of technological and economic feedbacks.
Davidson et al.'s matrix criterion concerns the propagation of feedback through
a coupled growth system. By contrast, \(\eta<\alpha\) rules out a particular
unbounded-value deviation in the developer's problem. Neither condition, by
itself, proves equilibrium existence.

I also relate the model to broader work on AI and economic growth.
\citet{aghionjonesjones2019} show how extending automation to idea production can
generate very rapid growth, while essential tasks that remain difficult to
automate can preserve bottlenecks; they use ``singularity'' for a finite-time
divergence. \citet{nordhaus2021} develops empirical diagnostics for an economic
singularity, while \citet{trammellkorinek2023} distinguishes the automation of
ordinary production from the automation of R\&D in a model of transformative AI.
\citet{jones2025rd} focuses directly on AI in research and development.
\citet{erdiletal2025gate} integrate a compute-based model of AI development, an
AI-automation block, and semi-endogenous growth with endogenous investment and
adjustment costs in a quantitative scenario model that abstracts from market
prices and market structure. Relative to these papers, I focus more narrowly on
how equilibrium prices and two factor-substitution margins govern the feedback,
the allocation of compute, and the distribution of output.

The terminology used to distinguish inference output, training activity,
algorithms, and model capital is informed by the economic accounting proposed by
\citet{korinekmckelvey2026}. Their framework treats compute as a central input
allocated between inference and training/R\&D and separates model capital from
algorithmic advances. In my reduced notation, raw inference compute and raw
research compute are paid uses of final output, while one capability stock
transforms both into economically useful services. Thus \(A\) combines objects
that their accounting keeps separate. This is a tractability assumption, not an
exact relabeling of their framework. I use their measurement distinctions; the
decentralized equilibrium and distributional results below do not follow from
that accounting framework.

The final-production block connects to the older literature on the elasticity of
substitution and growth \citep{pitchford1960,jonesmanuelli1990,rebelo1991}. CES
production is not a cosmetic generalization of Cobb--Douglas: whether the
elasticity lies below or above one governs whether labor remains a bottleneck or
AI services can asymptotically dominate along the paths studied here. Empirical work such as
\citet{duffypapageorgiou2000} illustrates why the aggregate elasticity should not
be imposed without scrutiny. \citet{aumshin2024} sharpen this point by estimating
complementarity between equipment and labor but gross substitution between
software and labor. Their estimates do not calibrate the model's AI-service elasticity,
but they support separating software-like inputs from undifferentiated physical
capital.
\citet{jonestonetti2026} provide a task-based counterpart to the model's
complementarity regime. With an elasticity of substitution below one across
tasks, the remaining tasks performed by slowly improving labor act as weak links
and delay the growth acceleration even when most tasks have been automated.
Their simulated aggregate paths under endogenous automation can remain difficult
to distinguish for a long interval before their growth rates diverge sharply. This is the closest
dynamic precedent for the delayed takeoff in my finite-boundary calculations.
An earlier illustration in \citet{aghionjonesjones2019} likewise shows that
bottlenecks can sustain relatively low aggregate growth while automation
continues, although their numerical automation path is exogenous.
Jones and Tonetti's simulations also show that automation alone does not determine the
limiting factor shares: the capital share can rise to one under complete
automation, fall toward zero when some weak links can never be automated, or
remain approximately stable along an intermediate path. The analogy is not an
identity. Their elasticity aggregates tasks and automation changes the set of
tasks assigned to capital; the model's \(\sigma_{XL}\) aggregates production labor and a
flow of quality-adjusted AI services.
The complementarity result adds an endogenous-research margin to this older CES
logic. \citet{davidsonetal2026} also show why, under CES aggregation, the relevant
object is the current elasticity of the appropriate aggregate with respect to its
AI input rather than a raw count of automated tasks. In my model,
\(s_X=\partial\ln Z/\partial\ln X\) and
\(s_M=\partial\ln E/\partial\ln(AM)\) play that role within
the two CES aggregators; the corresponding downstream elasticities are
\((1-\alpha)s_X\) for final output and \(\eta s_M\) for capability improvement.
Davidson et al. provide the closer analogue for feedback amplification and its
explosive-growth threshold, not for the endogenous timing of the low-growth
phase: their quantitative exercise imposes an automation change and fixes the
relevant research-resource shares. The simulated dates reported in those papers
are therefore not transferable to this economy.
These objects are endogenous CES shares rather than fractions of tasks. The adjustment is
also different: on the complementary-input branch I characterize, monopoly supply
keeps \(X/L\) finite and \(s_X\) positive while the raw-compute cost \(U=X/A\)
vanishes. On the regular complementary-input balanced-growth limit characterized
here, capability can continue to improve and research expenditure can become
almost fully automated while labor remains essential, output and consumption per
person remain bounded, and labor retains a positive income share.

The wage and factor-share results relate to the task and automation literature.
\citet{acemoglurestrepo2019} emphasize displacement and reinstatement of labor in
a task framework, and \citet{acemoglu2025} studies the macroeconomic gains from AI
through task exposure. That literature distinguishes tasks---units of work used
in production---from jobs, which bundle several tasks. Automation can therefore
change the tasks performed within a job without eliminating the job itself. My
model abstracts from task content, occupations, and the creation of new human
tasks. Labor is an aggregate input and full employment holds by construction.
The results should therefore be interpreted as model implications for wages,
factor shares, and sectoral labor allocation, not as predictions about the number
of jobs displaced. Within that abstraction, the model shows how a wage can rise
even when labor income as a share of output falls because the scale of output
changes endogenously.
This mechanism is related to the distinction between wages and labor income shares
emphasized by \citet{autorkausik2026}, but here it is generated by AI research
feedback and monopoly supply. \citet{trammellkorinek2023} likewise show that wages
can rise or fall when output expands and labor income as a share of output contracts. The conditional
AI-dominated limit I characterize maps that ambiguity into an explicit threshold for the
smaller of the substitution elasticities in final production and AI research.

\citet{restrepo2025agi} provides the closest antecedent to the limiting
distributional mechanism. When compute can perform all bottleneck work and
computational resources expand, output becomes linear in compute and labor,
growth is driven by compute, and labor's income share converges to zero. That
paper establishes the broad possibility that removing the human bottleneck can
make labor asymptotically inessential for aggregate growth. Taken together,
\citet{restrepo2025agi}, \citet{acemoglurestrepo2019}, and
\citet{trammellkorinek2023} establish both the displacement mechanism and the
possibility of a vanishing labor share. I claim neither as new.

My contribution is to place related mechanisms in a different
decentralized architecture and characterize their associated prices and
allocations under the hypotheses of each regime. The two CES elasticities
determine whether human labor remains a production or research bottleneck; the
integrated developer links the same capability stock to inference revenue and
research investment; and the consolidated accounts separate gross capital
payments, production and research wages, distributed developer profit, inference
compute, and research compute along the transition. This decomposition also
limits its own interpretation. Inference and research compute are resource uses,
not final income recipients. Without an upstream compute sector, the model cannot
determine whether the associated payments accrue to hardware producers, energy
suppliers, workers, or owners of additional capital.

Finally, the integrated developer connects the analysis to the emerging literature
on AI market structure \citep{korinekvipra2025,gans2024,demireretal2025}. Monopoly
and integration enter through separate channels: the developer charges a markup
on production AI services and internalizes the future operating profits created
by better capability. I hold market structure fixed in order to isolate the
technology. Comparing this
benchmark with monopolistic competition is a natural next step.
